\bookchapter{第二十二章：AI如何在手机中工作}

\chapterintro{欢迎来到手机最酷的“魔术现场”，主角是：人工智能。你每天在用的“一拍即得”的美照、张口就来的语音助手、越用越懂你的输入法，背后都站着这位隐形大脑。}

\booksection{一、AI相机：你的“后期暗房”搬进了手机}

现在的AI计算摄影，更像是把一个专业修图师+摄影师团队直接塞进了快门键里。在你按下快门前，AI就已经开始工作了：扫描场景，识别出“这是蓝天”、“这是人脸”、“这是夜景”。

\booksection{二、语音助手：你的“声音秘书”如何听懂你？}

\begin{enumerate}
  \item \textbf{听清（语音识别）}：将声波转变成文字。
  \item \textbf{听懂（自然语言理解）}：理解意图，提取关键实体。
  \item \textbf{执行与回答}：调用对应App或服务，再通过语音合成"说"给你听。
\end{enumerate}

\booksection{三、端侧AI vs. 云端AI}

端侧AI把计算留在手机里，速度快、隐私高、可离线。云端AI需要联网，但能处理更复杂的任务。

\booksection{四、AI大模型与手机：你的随身智能助手}

2025-2026年，AI大模型（如DeepSeek R1、ChatGPT-5、文心一言5.0、通义千问3等）的爆发深刻改变了手机的使用方式。手机不再只是执行指令的工具，而是能理解、推理和创作的智能伙伴。端侧AI芯片（NPU算力突破80 TOPS）让大模型推理首次在手机本地跑通，无需联网也能用AI。

\booksubsection{什么是AI大模型？}

AI大模型（Large Language Model，LLM）是通过海量文本数据训练出来的AI系统，能理解和生成自然语言。你可以把它想象成一个读过互联网上几乎所有文字的“超级学霸”，能回答问题、写文章、翻译、编程、分析数据。

\booksubsection{手机上的AI大模型应用：}

\begin{itemize}
  \item \textbf{对话助手}：DeepSeek、ChatGPT、豆包、Kimi等App，可以随时提问、翻译、写文案、总结文章。比传统搜索引擎更智能——它不只给你链接，而是直接给你答案。
  \item \textbf{写作与创作}：让AI帮你写邮件、总结会议纪要、润色文章、生成创意文案。
  \item \textbf{图像生成与编辑}：AI一键消除照片中的路人、扩展图片边界、风格转换、甚至根据文字描述生成图片。
  \item \textbf{实时翻译}：通话实时翻译、面对面对话翻译、拍照翻译，出国旅行不再有语言障碍。
  \item \textbf{学习辅导}：AI可以解释概念、出题练习、批改作业，是随身家教。
\end{itemize}

\booksubsection{端侧大模型：手机自己跑AI}

最新的旗舰手机已经能在本地运行小型AI大模型（通常7B参数左右），带来的变化是：

\begin{itemize}
  \item \textbf{不需要网络}：断网也能用AI功能，飞机上、地铁里都行。
  \item \textbf{隐私更安全}：你的对话和数据不会上传到云端，完全在手机本地处理。
  \item \textbf{响应更快}：没有网络延迟，AI回复几乎是即时的。
  \item \textbf{当前限制}：端侧模型的能力还比不上云端大模型，复杂任务仍需联网。但随着芯片AI算力的提升，端侧模型会越来越强。
\end{itemize}

\booksubsection{系统级AI集成：}

\begin{itemize}
  \item \textbf{苹果Apple Intelligence}：iOS 18引入的AI系统，集成写作工具、图像魔法棒、Siri增强、通知摘要等功能。端侧处理优先，复杂任务才上云端。
  \item \textbf{华为盘古大模型}：集成到鸿蒙系统中，小艺助手获得大模型能力，支持文档摘要、智慧识屏、AI修图等。
  \item \textbf{小米澎湃AI}：系统级AI能力，支持AI搜图、AI摘要、AI实时字幕、端侧大模型对话等。
  \item \textbf{三星Galaxy AI}：通话实时翻译、笔记AI摘要、图片生成编辑、搜索圈选等功能。
\end{itemize}

\booksection{五、几个典型应用场景与实操}

为了让你更直观地感受手机AI的“超能力”，下面我们用五个具体场景，展示AI是如何在日常生活中“出手相救”的：

\booksubsection{场景 1：实时翻译——出国旅行的随身翻译官}

\textbf{实际表现}：在巴黎点餐不会法语？iPhone的“翻译”App支持面对面双人对话模式，双方各看自己语言的字幕，甚至可以同声传译。安卓阵营的三星Galaxy AI通话实时翻译能让你在电话里直接说中文，对方听到的是英文，反之亦然。

\textbf{背后原理}：语音识别（ASR） 机器翻译（NMT） 语音合成（TTS）三步流水线，端到端延迟已可控制在1秒内。

\booksubsection{场景 2：AI消除——拯救被路人“毁掉”的游客照}

\textbf{实际表现}：在景点拍了一张绝美的自拍，但背景里挤满了其他游客？用“魔法消除”圈选路人，手机会自动识别并把路人“涂掉”，再智能填充背景，效果常常令人惊艳。

\textbf{背后原理}：图像分割（Inpainting）模型先识别并删除选中区域，然后根据周围像素和语义理解“脑补”出合理的填充内容。这背后的AI模型训练数据量高达数百万张图片。

\booksubsection{场景 3：AI修图——一键变“摄影大师”}

\textbf{实际表现}：拍了一张阴天的灰蒙蒙风景照？AI“魔法天空”可以一键把天空换成晴空、晚霞甚至星空。AI扩图可以让竖构图变成横构图，手机“脑补”出画面外的场景。

\textbf{背后原理}：基于生成式AI（Generative AI）的扩散模型（Diffusion Model），理解图像的语义内容后生成新像素。这是当前AI图像处理的“最前沿”。

\booksubsection{场景 4：文档处理——办公族的效率倍增器}

\textbf{实际表现}：会议结束后，把会议录音丢给AI，它能自动整理出结构化的会议纪要、提取待办事项、生成摘要。一份10页的英文PDF报告，AI可以在10秒内给你提炼出核心结论。

\textbf{背后原理}：长上下文大模型（Long Context LLM）可以“读完”数万字的文档，进行多轮推理和信息抽取。Apple Intelligence、华为小艺、WPS AI、腾讯元宝等都具备此能力。

\booksubsection{场景 5：AI搜图与识物——你的“视觉搜索引擎”}

\textbf{实际表现}：在街上看到一种不认识的花，相机一开，AI立刻告诉你“这是木槿花”。相册里翻找去年在海边拍的照片，搜索“大海”，AI自动把所有含海的照片找出来，不用手动翻。

\textbf{背后原理}：计算机视觉（CV）模型对图片内容进行向量化（Embedding），再与文字查询进行相似度匹配。这种“多模态搜索”是手机相册管理的一大飞跃。

\booksection{六、提示词优化：让 AI 听你指挥的“沟通公式”}

AI 的回答质量，很大程度上取决于你如何“提问”。同样一个问题，问法不同，AI 给的答案天差地别。下面我们介绍一套\textbf{最实用、最容易上手的“角色-背景-任务-限制”沟通思路}，并配以大量实例。

\begin{tcbinfo}{为什么需要框架？}
很多人用 AI 像用百度一样——敲几个关键词进去。但 AI 不是搜索引擎，它是一个“阅读理解后回答”的语言模型。\textbf{给它更多上下文、说清你要什么、限制它别跑偏}，它就能给你精准的结果。框架就是把你脑子里的“好答案长什么样”翻译成 AI 能理解的指令。

\end{tcbinfo}

\booksubsection{角色（Role）：告诉 AI 它是谁}

给 AI 一个具体的“身份”，它就会调用对应领域的专业知识和表达方式。\textbf{角色设定得越具体，输出越专业}。

\booksubsection{实用示例：}

\begin{itemize}
  \item "你是一位\textbf{有 10 年经验的小米手机维修工程师}。"（技术问题）
  \item "你是一位\textbf{擅长把复杂概念讲简单的中学物理老师}。"（科普问题）
  \item "你是一位\textbf{审美好、读得多、有自己风格的图书编辑}。"（写作问题）
  \item "你是一位\textbf{每天给上班族做早餐的烹饪博主}。"（生活问题）
  \item "你是一位\textbf{耐心、温和、永远不嘲笑新手的中医}。"（健康问题）
\end{itemize}

\textbf{为什么有效？}角色设定了 AI 的\textbf{知识范围}、\textbf{表达风格}和\textbf{专业深度}。没有角色的 AI 是个“全才但啥都不精”的通才；有了角色，它能“进入”那个身份，像那位专业人士一样思考。

\booksubsection{背景（Background）：告诉 AI 你的情况}

把“你是谁、你的现状、你的约束条件”提前说清楚，AI 就能给出\textbf{真正适合你}的答案，而不是泛泛而谈。

\booksubsection{实用示例：}

\begin{itemize}
  \item "我的手机是\textbf{小米 14，买了半年，最近玩游戏时发烫严重}。"（设备情况）
  \item "我是\textbf{一个 60 岁的退休阿姨，眼睛不太好，不太会用智能手机}。"（用户画像）
  \item "我正在\textbf{准备 2026 年 6 月的英语六级考试，基础一般}。"（学习背景）
  \item "我开了一家\textbf{三线城市的小咖啡店，月流水 3 万，雇了两个员工}。"（商业背景）
  \item "我\textbf{刚生完宝宝在家坐月子，睡眠严重不足}。"（生活阶段）
\end{itemize}

\textbf{为什么有效？}背景信息让 AI 知道\textbf{“为谁服务”}和\textbf{“在什么条件下”}。同一道题，给学生和给家长讲法不同；给城里人和给乡镇人建议也不同。背景缺失，AI 只能给“最大公约数”的答案，往往不痛不痒。

\booksubsection{任务（Task）：明确告诉 AI 做什么}

这一步要用\textbf{动词开头}（写、列、对比、翻译、解释、总结……），把任务描述得\textbf{具体、可执行}。含糊的任务只会换来含糊的答案。

\booksubsection{实用示例：}

\begin{itemize}
  \item "请\textbf{列出 5 个我应该立刻关闭的权限}，并说明原因。"（动作明确）
  \item "请\textbf{把下面这段英文邮件翻译成中文}，保持商务口吻。"（任务具体）
  \item "请\textbf{对比这四款 3000 元价位的手机}，给出一张优缺点对比表。"（格式明确）
  \item "请\textbf{用小学生能懂的话解释'什么是 5G'}。"（风格明确）
  \item "请\textbf{给我的早餐店写三条短视频文案}，分别针对上班族、情侣、带孩子的家长。"（细分明确）
\end{itemize}

\textbf{为什么有效？}任务部分是 AI 行动的“指令中枢”。\textbf{动词越清晰、要求越具体，AI 越知道“交付什么”}。模糊的“弄一下”、“搞个”、“看看”是 AI 大忌——它会揣测你的意思，结果常常跑偏。

\booksubsection{限制（Constraints）：告诉 AI 不要做什么}

“限制”听起来像“约束”，其实是\textbf{让 AI 收敛到最适合你的输出}的利器。常见维度有：字数、格式、语气、避免项、必含项。

\booksubsection{实用示例：}

\begin{itemize}
  \item "\textbf{不超过 200 字，分点列出}。"（长度+格式）
  \item "\textbf{语气要亲切自然，像朋友聊天}，不要用'首先、其次、再次'这种官话。"（语气）
  \item "\textbf{不要使用专业术语}，遇到专业名词要顺便解释。"（用词限制）
  \item "\textbf{每段不超过 3 行}，方便手机阅读。"（排版）
  \item "\textbf{必须包含：使用场景、优缺点、我的预算匹配度}。"（必含项）
  \item "\textbf{不要给具体投资建议}，只解释原理。"（边界）
\end{itemize}

\textbf{为什么有效？}没有限制，AI 容易“自由发挥”——要么太学术、要么太长、要么跑题。明确的限制=明确的“交付标准”，双方对齐预期，结果才可控。

\booksubsection{提示词质量评估：好提示词的 5 个特征}

写好提示词不是一次性的事，\textbf{写完先自查这五条}，再发给 AI：

\begin{enumerate}
  \item \textbf{具体而非抽象}：不是"写篇文章"，而是"写一篇 800 字、3 个小节、面向新手妈妈的手机选购指南"。
  \item \textbf{包含背景而非孤立}：告诉 AI"我是谁、为什么问这个"，它才能给出"为你定制"的答案。
  \item \textbf{有明确输出格式}：是表格、列表、段落、还是 JSON？提前说。
  \item \textbf{限制到位}：字数、风格、避项，三选一就有明显提升。
  \item \textbf{可被 AI 验证}：问完之后，让 AI 复述一遍你的要求；说不清楚就是没说清楚。
\end{enumerate}

\booksubsection{常见问题与解法}

\begin{center}
\begin{tabularx}{\linewidth}{|X|X|}
\hline
\textbf{症状} & \textbf{解法} \tabularnewline
\hline
回答太泛、不实用 & 加背景（你是谁/什么场景） \tabularnewline
回答太学术/太官方 & 加语气限制 + 拒绝"官话"清单 \tabularnewline
回答跑题/答非所问 & 把任务用动词重新说一遍，并要求"只回答这个" \tabularnewline
回答太长看不完 & 硬性限制字数（200 字内） \tabularnewline
AI 一本正经胡说 & 加"如不确定请标注'我不确定'" + 重要事实自查 \tabularnewline
第一版不满意 & 不要重写，在原提示词后追加"再…一些" \tabularnewline
\hline
\end{tabularx}
\end{center}

\booksubsection{三个完整场景案例}

\begin{tcbinfo}{场景一：生活助手（早餐搭配建议）}
\textbf{提示词}：

你是一位擅长营养搭配的早餐博主（\textbf{角色}）。我是一位上班族妈妈，每天早上只有 15 分钟做早餐，家里有 6 岁的孩子和 60 岁的婆婆（\textbf{背景}）。请\textbf{列出 7 天不重样的早餐搭配}（\textbf{任务}）。每份早餐要求：(1) 准备时间不超过 15 分钟 (2) 营养均衡（蛋白+碳水+蔬果）(3) 老少皆宜 (4) 食材在普通超市能买到（\textbf{限制}）。

\textbf{为什么有效}：角色定专业度，背景给三类人共性，任务明确交付物，4 条限制确保可执行。

\end{tcbinfo}

\begin{tcbinfo}{场景二：学习辅助（英语六级备考）}
\textbf{提示词}：

你是一位\textbf{经验丰富的英语六级老师}（角色）。我是一名普通本科生，\textbf{四级 480 分，六级目标 500，还有 3 个月}（背景）。请\textbf{为我制定一份可执行的 12 周备考计划}（任务）。要求：(1) 每周明确听说读写的时间分配 (2) 给出每个月的重点突破方向 (3) 推荐 3 本合适的复习资料（中文说明理由） (4) 表格形式输出（限制）。

\textbf{为什么有效}：角色定专业背景，背景给出分数和时间锚点，任务明确“分阶段计划”，限制要求“表格化”和“可执行”。

\end{tcbinfo}

\begin{tcbinfo}{场景三：工作效率（会议纪要整理）}
\textbf{提示词}：

你是一位\textbf{资深行政助理}（角色）。我刚开完一个 1 小时的部门会议，\textbf{讨论 Q2 营销推广方案}（背景）。\textbf{请把以下会议记录整理成结构化纪要}（任务）。要求：(1) 分“决议事项 / 待办事项 / 风险与异议”三部分 (2) 待办事项标出负责人和截止时间（如未提到则写“待确认”） (3) 总字数控制在 600 字内 (4) 用 markdown 格式输出（限制）。

\textbf{为什么有效}：角色定行政视角，背景给会议主题，任务明确结构化输出，4 条限制让结果直接可复用。

\end{tcbinfo}

\textbf{进阶玩法：让 AI 成为你的“私人助理”}\par

试试这些进阶用法，体验 AI 的真正实力：

\begin{itemize}
  \item \textbf{多轮对话优化}：对 AI 的回答不满意时，可以说"这个太专业了，请用小学生能懂的话再解释一次"。
  \item \textbf{链式思考（CoT）}：让 AI"一步步思考"，回答质量通常更高："请一步步分析这道数学题"。
  \item \textbf{多 AI 协作}：用 A 写初稿，B 润色，C 翻译——组合使用效果更佳。
  \item \textbf{系统提示词}：在 AI 设置中预设你的身份、偏好、忌讳等，让回答更贴合你的需求。
\end{itemize}

\booksection{七、AI的局限与挑战：保持清醒的判断力}

AI虽然强大，但绝非“无所不能”。清醒地认识它的局限，才能更好地使用它：

\begin{itemize}
  \item \textbf{"幻觉"问题}：AI可能一本正经地胡说八道（专业术语叫Hallucination）。它会编造不存在的历史事件、错误的引言、虚假的统计数据。重要信息必须人工核实。
  \item \textbf{知识截止}：大多数AI的训练数据有截止日期（如2024年初），它不知道之后发生的事。
  \item \textbf{推理能力有限}：复杂的多步推理、数学证明、逻辑陷阱，AI仍然容易出错。
  \item \textbf{隐私风险}：在云端AI服务中，你的对话可能被用于模型训练。敏感信息应优先使用端侧AI。
  \item \textbf{偏见与价值观}：AI可能带有训练数据中的偏见，需要保持独立思考。
\end{itemize}

\begin{tcbinfo}{AI使用小贴士}
\begin{itemize}
  \item AI生成的内容可能不准确（"幻觉"问题），重要信息务必核实。
  \item 不要向AI输入个人隐私信息（身份证号、密码、银行信息等）。
  \item 善用提示词（Prompt）：越具体、越清晰的描述，AI的回答质量越高。
  \item AI是工具不是权威，保持独立思考和判断力。
  \item 遇到关键决策（医疗、法律、财务），咨询专业人士而非依赖AI。
\end{itemize}

\end{tcbinfo}

\begin{tcbpractice}{来挑战：三道AI实战题}
\begin{enumerate}
  \item \textbf{挑战一：复杂指令}：对语音助手说："明早7点提醒我带伞，并把今天下午3点的日历改到4点。" 看看它能否一次完成多步操作。
  \item \textbf{挑战二：AI修图}：拍一张有路人的游客照，使用"魔法消除"或"AI橡皮擦"功能，把路人清除干净。观察AI对背景的"脑补"效果。
  \item \textbf{挑战三：提示词优化}：先用简单提示词"写一段旅游攻略"问AI，再用CRISPE框架写一个详细提示词。比较两次回答的质量差异。
\end{enumerate}

\end{tcbpractice}

\begin{tcbinfo}{案例 2：换机数据迁移完整流程}
\begin{itemize}
  \item \textbf{场景}：iPhone 13 小米 15 跨品牌换机。
  \item \textbf{步骤}：(1) 旧机 iCloud/微信自带备份 (2) 退出 Apple ID + 关闭"查找" (3) 新机用"小米换机"App 扫码 (4) 选择迁移项 (5) 微信单独用自带"聊天记录迁移" (6) 验证关键数据（通讯录/相册/微信） (7) 旧机加密擦除。
  \item \textbf{预期结果}：数据零丢失，旧机可安全转卖。
\end{itemize}

\end{tcbinfo}
